\documentclass[10pt,a4paper]{amsart}
\usepackage{amsmath,amssymb,amsthm,graphicx}
% Retained preamble for the Stacks Project chapter snapshot in this directory.
% The chapter's own preamble.tex is a copy of the shared Stacks Project preamble; the block
% below is sliced MECHANICALLY, line for line, from it. Nothing was added, removed or
% redefined; the wrapper class, the xr-hyper/hyperref wiring and the externaldocument list
% are omitted because the assembler supplies the document class.
%
% lines taken from preamble.tex: 3, 5, 9, 11, 12, 156, 157, 158, ... (51 lines total)
%
% package coverage restored from the chapter's own preamble (2026-10-13): xy, xy-2cell

\usepackage[all]{xy}
\xyoption{2cell}
\UseAllTwocells

\newenvironment{reference}{\comment}{\endcomment}
\newenvironment{slogan}{\comment}{\endcomment}
\newenvironment{history}{\comment}{\endcomment}
\theoremstyle{plain}
\newtheorem{theorem}[subsection]{Theorem}
\newtheorem{proposition}[subsection]{Proposition}
\newtheorem{lemma}[subsection]{Lemma}
\theoremstyle{definition}
\newtheorem{definition}[subsection]{Definition}
\newtheorem{example}[subsection]{Example}
\newtheorem{exercise}[subsection]{Exercise}
\newtheorem{situation}[subsection]{Situation}
\theoremstyle{remark}
\newtheorem{remark}[subsection]{Remark}
\newtheorem{remarks}[subsection]{Remarks}
\numberwithin{equation}{subsection}
\def\lim{\mathop{\mathrm{lim}}\nolimits}
\def\colim{\mathop{\mathrm{colim}}\nolimits}
\def\Spec{\mathop{\mathrm{Spec}}}
\def\Hom{\mathop{\mathrm{Hom}}\nolimits}
\def\Ext{\mathop{\mathrm{Ext}}\nolimits}
\def\SheafHom{\mathop{\mathcal{H}\!\mathit{om}}\nolimits}
\def\SheafExt{\mathop{\mathcal{E}\!\mathit{xt}}\nolimits}
\def\Sch{\mathit{Sch}}
\def\Mor{\mathop{\mathrm{Mor}}\nolimits}
\def\Ob{\mathop{\mathrm{Ob}}\nolimits}
\def\Sh{\mathop{\mathit{Sh}}\nolimits}
\def\NL{\mathop{N\!L}\nolimits}
\def\CH{\mathop{\mathrm{CH}}\nolimits}
\def\proetale{{pro\text{-}\acute{e}tale}}
\def\etale{{\acute{e}tale}}
\def\QCoh{\mathit{QCoh}}
\def\Ker{\mathop{\mathrm{Ker}}}
\def\Im{\mathop{\mathrm{Im}}}
\def\Coker{\mathop{\mathrm{Coker}}}
\def\Coim{\mathop{\mathrm{Coim}}}
\DeclareMathSymbol{\boxtimes}{\mathbin}{AMSa}{"02}
\def\QCohstack{\mathcal{QC}\!\mathit{oh}}
\def\Cohstack{\mathcal{C}\!\mathit{oh}}
\def\Spacesstack{\mathcal{S}\!\mathit{paces}}
\def\Quotfunctor{\mathrm{Quot}}
\def\Hilbfunctor{\mathrm{Hilb}}
\def\Curvesstack{\mathcal{C}\!\mathit{urves}}
\def\Polarizedstack{\mathcal{P}\!\mathit{olarized}}
\def\Complexesstack{\mathcal{C}\!\mathit{omplexes}}
\def\Pic{\mathop{\mathrm{Pic}}\nolimits}
\def\Picardstack{\mathcal{P}\!\mathit{ic}}
\def\Picardfunctor{\mathrm{Pic}}
\def\Deformationcategory{\mathcal{D}\!\mathit{ef}}
\usepackage{cleveref}
\crefname{originalequation}{original equation}{original equations}
\Crefname{originalequation}{Original equation}{Original equations}
\makeatletter
\newlabel{equation-pi}{{08PR}{source}{}{sourcepointer-1}{}}
\newlabel{equation-pi@cref}{{[originalequation][0][]08PR}{[1][0][]source}{}{}{}}
\makeatother
\crefname{originallemma}{original lemma}{original lemmas}
\Crefname{originallemma}{Original lemma}{Original lemmas}
\makeatletter
\newlabel{lemma-identify-pi-shriek}{{08PS}{source}{}{sourcepointer-2}{}}
\newlabel{lemma-identify-pi-shriek@cref}{{[originallemma][0][]08PS}{[1][0][]source}{}{}{}}
\makeatother
\crefname{originallemma}{original lemma}{original lemmas}
\Crefname{originallemma}{Original lemma}{Original lemmas}
\makeatletter
\newlabel{simplicial-lemma-homotopy-equivalence}{{08P2}{source}{}{sourcepointer-3}{}}
\newlabel{simplicial-lemma-homotopy-equivalence@cref}{{[originallemma][0][]08P2}{[1][0][]source}{}{}{}}
\makeatother
\crefname{originallemma}{original lemma}{original lemmas}
\Crefname{originallemma}{Original lemma}{Original lemmas}
\makeatletter
\newlabel{simplicial-lemma-qis-simplicial-abelian-groups}{{08P1}{source}{}{sourcepointer-4}{}}
\newlabel{simplicial-lemma-qis-simplicial-abelian-groups@cref}{{[originallemma][0][]08P1}{[1][0][]source}{}{}{}}
\makeatother
\crefname{originallemma}{original lemma}{original lemmas}
\Crefname{originallemma}{Original lemma}{Original lemmas}
\makeatletter
\newlabel{simplicial-lemma-standard-simplicial-homotopy}{{08ND}{source}{}{sourcepointer-5}{}}
\newlabel{simplicial-lemma-standard-simplicial-homotopy@cref}{{[originallemma][0][]08ND}{[1][0][]source}{}{}{}}
\makeatother
\title{The lower shriek computed by the standard resolution}
\author{The Stacks Project Authors}
\date{stacks-project:cotangent.tex:local-snapshot-20261012; GNU FDL 1.2 or later}
\begin{document}
\maketitle
\setcounter{subsection}{1}
\noindent\textit{Editorial note.} All delivered mathematical text is the author's own, copied from the pinned Stacks Project chapter file cotangent.tex; no statement or proof step was rewritten, reordered or removed.\par\medskip
\noindent\textit{Editorial note.} The proof environment carries the author's optional label 'First proof'; the optional argument is part of the counted proof span and is preserved verbatim.\par\medskip
\noindent\textit{Editorial note.} The five references of the counted spans are recorded as non-load-bearing navigation pointers; the argument that consumes them stays inside the counted proof span.\par\medskip
\noindent\textit{Editorial note.} Preamble: the retained preamble of this chapter (sources/stacks-cotangent/retained-preamble.tex, lines 1-63), containing the Xy-pic setup, the comment environments, the theorem declarations and the macros including \textbackslash{}NL.\par\medskip
\noindent\textit{Allowed standard prerequisites (editorial).}\begin{itemize}
\item The identification of the lower shriek by polynomial resolutions with a trivial Kan fibration. Lemma lemma-identify-pi-shriek (445-456, produced as L44140), used at proof lines 506-510. Reduces the claim to its two hypotheses.
\item The standard resolution map is a homotopy equivalence. Simplicial, Lemma simplicial-lemma-standard-simplicial-homotopy, used at proof lines 510-513. Gives the first step of the chain.
\item A homotopy equivalence induces a quasi-isomorphism of associated complexes. Simplicial, Lemma simplicial-lemma-homotopy-equivalence, used at proof lines 513-516. Gives the second step.
\item A quasi-isomorphism of associated complexes makes the map a trivial Kan fibration. Simplicial, Lemma simplicial-lemma-qis-simplicial-abelian-groups, used at proof lines 516-519. Concludes the second hypothesis.
\end{itemize}\medskip
\ifdefined\hypertarget\hypertarget{sourcepointer-1}{}\fi
\noindent\textit{Source reference.} Original equation 08PR: focused-seeds/AGP/sources/stacks-snapshot-20261011/tags.txt, source lines 11034--11034. Navigation pointer into the Stacks Project; the standard fact it supplies is applied inside the counted proof span and the argument resting on it stays there.\par\medskip
\ifdefined\hypertarget\hypertarget{sourcepointer-2}{}\fi
\noindent\textit{Source reference.} Original lemma 08PS: focused-seeds/AGP/sources/stacks-snapshot-20261011/tags.txt, source lines 11035--11035. Navigation pointer into the Stacks Project; the standard fact it supplies is applied inside the counted proof span and the argument resting on it stays there.\par\medskip
\ifdefined\hypertarget\hypertarget{sourcepointer-3}{}\fi
\noindent\textit{Source reference.} Original lemma 08P2: focused-seeds/AGP/sources/stacks-snapshot-20261011/tags.txt, source lines 11010--11010. Navigation pointer into the Stacks Project; the standard fact it supplies is applied inside the counted proof span and the argument resting on it stays there.\par\medskip
\ifdefined\hypertarget\hypertarget{sourcepointer-4}{}\fi
\noindent\textit{Source reference.} Original lemma 08P1: focused-seeds/AGP/sources/stacks-snapshot-20261011/tags.txt, source lines 11009--11009. Navigation pointer into the Stacks Project; the standard fact it supplies is applied inside the counted proof span and the argument resting on it stays there.\par\medskip
\ifdefined\hypertarget\hypertarget{sourcepointer-5}{}\fi
\noindent\textit{Source reference.} Original lemma 08ND: focused-seeds/AGP/sources/stacks-snapshot-20261011/tags.txt, source lines 10986--10986. Navigation pointer into the Stacks Project; the standard fact it supplies is applied inside the counted proof span and the argument resting on it stays there.\par\medskip
\begin{lemma}
\label{lemma-pi-shriek-standard}
Let $A \to B$ be a ring map. Let $\epsilon : P_\bullet \to B$
be the standard resolution of $B$ over $A$. Let $\pi$ be as in
(\ref{equation-pi}). Then
$$
L\pi_!(\mathcal{F}) = \mathcal{F}(P_\bullet, \epsilon)
$$
in $D(\textit{Ab})$, resp.\ $D(B)$ functorially in $\mathcal{F}$ in
$\textit{Ab}(\mathcal{C})$, resp.\ $\textit{Mod}(\underline{B})$.
\end{lemma}

\begin{proof}[First proof]
We will apply Lemma \ref{lemma-identify-pi-shriek}.
Since the terms $P_n$ are polynomial algebras we see the first
assumption of that lemma is satisfied. The second assumption is proved
as follows. By
Simplicial, Lemma \ref{simplicial-lemma-standard-simplicial-homotopy}
the map $\epsilon$ is a homotopy equivalence of underlying
simplicial sets. By
Simplicial, Lemma \ref{simplicial-lemma-homotopy-equivalence}
this implies $\epsilon$ induces a quasi-isomorphism of associated
complexes of abelian groups. By
Simplicial, Lemma \ref{simplicial-lemma-qis-simplicial-abelian-groups}
this implies that $\epsilon$ is a trivial Kan fibration of underlying
simplicial sets.
\end{proof}

\end{document}
